\chapter{Resumen del Análisis}

En este anexo, el análisis realizado demuestra la viabilidad técnica y funcional del sistema de gestión financiera personal propuesto. 
La combinación de tecnologías modernas como OCR para categorización automática, y metodologías financieras probadas como el sistema Dave Ramsey, 
proporciona una base sólida para el desarrollo.

\section{Principales Hallazgos}

\subsection{Fortalezas del Diseño}
\begin{itemize}
    \item La automatización del registro de gastos mediante OCR reduce significativamente la fricción para el usuario
    \item La limitación a 10 categorías fomenta la simplicidad y evita la sobre-categorización
    \item La arquitectura modular permite escalabilidad y mantenimiento eficientes
    \item El enfoque en jóvenes adultos aborda un nicho específico con necesidades claras
\end{itemize}

\subsection{Consideraciones Técnicas}
\begin{itemize}
    \item La precisión del OCR depende de la calidad de las imágenes capturadas
    \item Se requiere conectividad para funcionalidades en la nube, pero se mantiene funcionalidad básica offline
    \item La categorización automática mejorará con el uso y recopilación de datos
\end{itemize}

\section{Impacto Esperado}

\subsection{Para los Usuarios}
\begin{itemize}
    \item Reducción del tiempo dedicado al registro manual de gastos
    \item Mayor conciencia sobre patrones de gasto
    \item Desarrollo de hábitos financieros saludables
    \item Mejor planificación y control de deudas
\end{itemize}

\subsection{Para el Ámbito Académico}
\begin{itemize}
    \item Demostración práctica de integración de IA basica y tecnologías como el OCR en aplicaciones financieras
    \item Validación de metodologías de educación financiera en formato digital
    \item Contribución al desarrollo de fintech en México
\end{itemize}